\chapter*{Список использованных сокращений}
\addcontentsline{toc}{chapter}{Список использованных сокращений}
\begin{description}
\item[EMG] Электромиограмма (electromyogram).
\item[ECAP] Вызванный составной потенциал действия (evoked compound action potential).
\item[FEM] Метод конечных элементов (finite element method).
\item[LIF] Нейронная модель с утечкой, интегрированием и пороговым разрядом (leaky integrate-and-fire).
\item[RMSE] Корень из средней квадратичной ошибки (root mean squared error).
\item[SCS] Стимуляция спинного мозга (spinal cord stimulation).
\end{description}

\chapter*{Список терминов}
\addcontentsline{toc}{chapter}{Список терминов}
\noindent динамическое представление: Формализованное описание выбранных временных характеристик нейронной активности.\par
\noindent коннектом: Совокупность сведений о нейронах и связях между ними.\par
\noindent кривая роста ECAP: Зависимость амплитуды вызванного составного потенциала действия от амплитуды стимуляции при фиксированных остальных параметрах.\par
\noindent оператор наблюдения: Преобразование состояния модели в величину, сопоставимую с измерением.\par
\noindent рекрутирование волокон: Вовлечение нервных волокон в вызванный стимуляцией ответ.\par
